\documentclass[letterpaper,11pt]{article}
\usepackage[top=0.8in, bottom=0.8in, left=0.9in, right=0.9in]{geometry}
\usepackage[hidelinks]{hyperref}
\usepackage{lmodern}
\usepackage{parskip}
\usepackage{microtype}
\setlength{\parskip}{6pt}
\setlength{\parindent}{0pt}
\begin{document}
\begin{flushleft}
\textbf{\large Ajay Venkatesh} \\
Brooklyn, NY \\
+1 516 585 9013 \\
\href{mailto:av3855@nyu.edu}{av3855@nyu.edu}
\end{flushleft}
\vspace{10pt}
May 21, 2026
\vspace{6pt}

Hiring Manager \\
Wolverine Trading

\vspace{10pt}

Dear Hiring Manager,

My name is Ajay Venkatesh. I'm finishing my M.S. in Computer Engineering at NYU, with production software experience at PwC in Bengaluru, a summer SDE internship at Amazon, and ongoing on-campus work at NYU's Novel AI Technologies. I'm writing about the Entry-Level C++ Software Engineer role at Wolverine Trading.

My graduate coursework has been weighted toward the systems end. \textbf{Computer Architecture}, \textbf{System Design}, \textbf{Data Structures}, \textbf{Big Data}, and \textbf{Advanced Java} alongside \textbf{Machine Learning} and \textbf{Deep Learning}. \textbf{C++} is where I cut my teeth and where I keep going back when an abstraction stops being honest about what's happening underneath. Object-oriented design, multithreaded programming, and performance tuning are the parts of engineering I find most interesting, especially when the runtime matters as much as the correctness.

On engineering rigor: at Amazon I shipped backend services in Java (Coral, Smithy, Dagger, CBOR) with low-level design patterns and unit / integration testing via \textbf{Mockito}, contributed under code-review discipline, and treated tests as part of the build. I provisioned \textbf{Infrastructure-as-Code} on \textbf{AWS CDK} for distributed message processing at high throughput. The discipline of writing clean, profilable, and well-tested code carries over directly to a C++ environment where every instruction matters.

On data at scale: my Big Data project built a \textbf{PySpark / HDFS} pipeline over millions of flight records, training Random Forest and Gradient Boosted Tree models with rigorous EDA. The shape of the problem, feature pipelines, anomaly detection, and squeezing signal out of irregular real-world data, is adjacent to quant work even at a much smaller scale.

What pulls me toward Wolverine is the discipline. Building trading systems where the work is measured in real units, on real markets, alongside people who think probabilistically about risk and reward is the kind of engineering culture I want to grow inside of. I'm at the start of my career and looking for a place where the bar is high.

I'd welcome the chance to discuss further. Thanks for considering my application.

\vspace{12pt}
Sincerely, \\
Ajay Venkatesh

\end{document}